\section{Logical target and standards of formalization}

\subsection{iPRA as used in the repository}

For this report, iPRA is first-order intuitionistic logic with equality over a
function symbol for every primitive-recursive construction, the defining
equations for those symbols, the usual successor separation and injectivity
axioms, and induction for quantifier-free formulas.  The checked repository
implementation makes this precise:
\begin{itemize}
\item \(\PR(k)\) is the inductive type of primitive-recursive symbols of arity
\(k\), generated by zero, successor, projections, composition, and primitive
recursion;
\item \(\PTm(n)\) is the type of terms with at most \(n\) free de Bruijn
variables;
\item \(\PRAAx\) is the family of open iPRA axioms;
\item \(\Prf\,\PRAAx\,\Gamma\,\varphi\) is intuitionistic derivability from
context \(\Gamma\).
\end{itemize}

The restriction to quantifier-free induction is not a serious obstacle to
finite computation.  It is a design constraint: state, loop counters,
certificates, and all bounded quantifiers must be compiled into terms and
Boolean equations before induction is invoked.

\subsection{Three levels that must not be conflated}

A robust formalization keeps three levels visible.

\paragraph{Meta-level specification.}
Lean structures may organize arities, formulas, and proof terms.  For example,
a structure may contain fixed symbols \(\Good,\Run,\Inv,\Accept\) and fields
whose types are local \(\Prf\,\PRAAx\) derivations.  This is harmless: Lean is
only assembling an explicit object-language proof.

\paragraph{Host-level semantic validation.}
An executable program can test examples, discover missing bounds, and validate
finite certificate formats.  Such tests establish that the proposed finite
algorithm behaves as intended in the standard model.  They do not establish
that iPRA proves its correctness.

\paragraph{Object-language derivation.}
The final theorem is a constructor returning \(\Prf\,\PRAAx\).  Every use of
induction, equality transport, bounded search, or a local theorem is visible in
that constructor.  Only this level is transported by \code{pra_iff_native}.

\begin{warning}[Semantic reflection is not enough]
A theorem of the form
\[
  \forall u\in\N,\quad \operatorname{run}(u)\text{ satisfies }\Accept
\]
proved in Lean's ambient logic is not interchangeable with
\[
  []\vdash[\PRAAx]\forall u\,\Accept(\Run(\operatorname{bound}(u),u),u).
\]
The former may use unrestricted induction, higher-order choice, or library
facts unavailable to iPRA.  The framework requires local object-language
premises and constructs the latter directly.
\end{warning}

\subsection{Fixed families rather than a universal evaluator}

There is no total primitive-recursive universal evaluator for all primitive-
recursive functions.  Therefore a valid framework cannot quantify over an
arbitrary code \(e\) and assume a total PR term \(U(e,x)\) evaluating it.
There are three safe alternatives.

\begin{enumerate}
\item \textbf{Fixed symbols.}  For each formalized theorem, choose finitely many
symbols in \(\PR(k)\).  This is the simplest and strongest route.
\item \textbf{Fixed program family.}  Define an inductive family of programs
whose interpreter is primitive recursive because a structural size or a
syntactic recursion depth is fixed by the family.
\item \textbf{Clocked computation.}  Evaluate a code for at most \(t\) steps.
The bounded evaluator is primitive recursive; a trace certificate supplies
\(t\) and proves acceptance.
\end{enumerate}

The framework is parametric at the Lean level but concrete in the object
language.  A theorem may accept a structure containing arbitrary \emph{fixed}
\(\PR\) symbols.  It never asserts the existence of one internal evaluator for
all such structures.

\subsection{The target shape of a reusable theorem}

The basic compiler theorem has the following logical form.  Let
\(G(u)\), \(I(i,s,u)\), and \(A(s,u)\) be quantifier-free predicates represented
by Boolean equations, and let \(r(i,u)\) be the primitive-recursive run.
Assume closed derivations of
\begin{align*}
 &G(u)\to I(0,r(0,u),u), \\
 &G(u)\to I(i,s,u)\to I(Si,\operatorname{step}(i,s,u),u),\\
 &G(u)\to I(b(u),s,u)\to A(s,u).
\end{align*}
Then the compiler must return a closed derivation of
\[
   \forall u\,[G(u)\to A(r(b(u),u),u)].
\]
The theorem is deliberately intensional.  Its inputs are not claims that the
three implications are true in \(\N\); they are values of the corresponding
\(\Prf\,\PRAAx\) types.

\subsection{What counts as a complete Lean outline}

A useful outline must expose enough information that completion is routine
proof engineering rather than renewed mathematical research.  In this package
that means:
\begin{itemize}
\item every term builder has a fixed arity and a documented de Bruijn layout;
\item every formula used by induction is explicitly defined;
\item the quantifier-free side condition is a named lemma;
\item local hypotheses and the global result have exact
\(\Prf\,\PRAAx\) types;
\item the zero and successor run equations are tied directly to
\code{PRA.precZero} and \code{PRA.precSucc};
\item the remaining proof is split into base, step, induction, and final
milestones, each using named checked APIs;
\item every unproved part remains marked by \code{sorry}; no compilation claim
is made.
\end{itemize}

This is stronger than a pseudocode sketch but weaker than a completed kernel
proof.  It gives a clear path to completion and makes the residual work
measurable.
